\chapter{Grands modèles de langage en finance}\label{ch:ml:large-language-models-in-finance}

Un \omterm{def:ml:large-language-models-in-finance:lm}{modèle de langage} à qui l'on demande si des titres de sa propre période d'entraînement sont une bonne nouvelle pour l'action a raison dans 70~\% des
cas~; interrogé sur des titres rédigés après la fin de ses données d'entraînement, il a raison dans 56~\% des cas. Le premier chiffre mesure
sa mémoire, le second son jugement, et un backtest qui mélange les deux ne mesure rien. Le modèle en question est celui de ce
chapitre~: un \omterm{def:ml:sequence-models-on-order-books:attention}{transformer} de 107\,000 paramètres entraîné en une minute sur les titres synthétiques du chapitre~13, assez petit pour
être réentraîné avec et sans la période d'évaluation dans ses données. Le protocole qu'il teste est celui dont un desk a besoin pour
les modèles commerciaux dont il ne voit pas les données d'entraînement. Le chapitre passe ensuite à l'usage des \omterm{def:ml:large-language-models-in-finance:lm}{modèles de langage} en
pratique~: lecture à grande échelle de dépôts et d'actualités par recherche documentaire, exécution de tâches de recherche comme agents, et ce qu'un desk peut
et ce qu'il ne peut pas leur confier. Rien dans ce chapitre n'appelle un service en réseau et ne télécharge de modèle~; chaque chiffre provient de code local
accompagné de tests.

\section{Modèles de langage et leur entraînement}

\begin{definition}[Modèle de langage, grand modèle de langage]\label{def:ml:large-language-models-in-finance:lm}
Un \emph{modèle de langage}\index{modèle de langage} attribue des probabilités à des séquences de jetons, généralement par la probabilité conditionnelle
de chaque jeton sachant les précédents, $\P(w_t\mid w_1,\dots,w_{t-1})$~; il génère un texte en tirant un
jeton à la fois. Un \emph{grand modèle de langage}\index{grand modèle de langage} (LLM) est un modèle de langage à base de \omterm{def:ml:sequence-models-on-order-books:attention}{transformer} (chapitre~8)
à plusieurs milliards de paramètres, préentraîné (chapitre~10) sur une large part des textes publics d'internet et des livres, puis
affiné pour suivre des instructions.
\end{definition}

\begin{definition}[Prompt, fenêtre de contexte, date de coupure de l'entraînement]\label{def:ml:large-language-models-in-finance:prompt}
Un \emph{prompt}\index{prompt} est le texte donné à un \omterm{def:ml:large-language-models-in-finance:lm}{modèle de langage}, qu'il poursuit. La \emph{fenêtre de
contexte}\index{fenêtre de contexte} est le plus grand nombre de jetons sur lesquels le modèle peut conditionner à la fois, prompt et réponse
réunis. La \emph{date de coupure de l'entraînement}\index{date de coupure de l'entraînement} est la date du texte le plus récent des données d'entraînement du modèle~:
le modèle peut connaître tout ce qui a été publié avant cette date et ignore tout ce qui l'a été après, sauf ce que lui indique le prompt.
\end{definition}

Le \Cref{lst:ml:llm:tinylm} présente l'architecture à échelle réduite~: plongements de jetons et de positions, deux couches
d'auto-attention causale (un jeton ne porte son attention que sur les jetons précédents), et une projection vers les 505 mots du vocabulaire. Il est
entraîné, comme ses grands cousins, à prédire chaque jeton suivant. Son texte d'entraînement est le corpus du chapitre~13, tronqué à ses premiers
10\,000 titres, écrit sous forme de phrases telles que \texttt{w190 apex mining beat estimates amid copper demand => up}~: une semaine,
l'entreprise, le titre, et, après la flèche, ce qu'a fait l'action (hausse ou baisse, soit le signe du rendement de réaction). Interroger le
modèle sur un titre revient à lui soumettre tout ce qui précède la flèche et à comparer les probabilités qu'il attribue à
\texttt{up} et \texttt{down}. Deux exemplaires sont entraînés~: un modèle \emph{propre}, dont les données d'entraînement s'arrêtent au jour 907, et un
modèle \emph{contaminé}, dont les données vont jusqu'au jour 1\,210~; les deux sont évalués sur les 5\,972 titres antérieurs au jour 907, les 2\,016
titres compris entre les deux dates de coupure et les 2\,012 postérieurs aux deux.

Les grands modèles se construisent de la même façon, à une échelle que ce chapitre ne peut atteindre~: BloombergGPT (Wu et coauteurs, 2023), par
exemple, est un modèle de 50 milliards de paramètres entraîné sur un jeu de données financières de 363 milliards de jetons, construit à partir des sources de Bloomberg
auxquelles s'ajoutent 345 milliards de jetons de texte général. Ce qui se transpose du modèle réduit, c'est le protocole, et non les chiffres.

\section{Dépôts, conférences et actualités à grande échelle~: recherche documentaire}

Un \omterm{def:ml:large-language-models-in-finance:lm}{modèle de langage} ne connaît que ses données d'entraînement et son \omterm{def:ml:large-language-models-in-finance:prompt}{prompt}. Pour répondre à des questions sur une collection de documents qui change
chaque jour (dépôts, transcriptions, actualités), on lui fournit les passages pertinents dans son \omterm{def:ml:large-language-models-in-finance:prompt}{prompt}.

\begin{definition}[Génération augmentée par la recherche, BM25]\label{def:ml:large-language-models-in-finance:rag}
La \emph{génération augmentée par la recherche}\index{génération augmentée par la recherche} (RAG) répond à une question en récupérant
d'abord les passages pertinents pour y répondre dans une collection de documents, puis à produire la réponse à partir d'un \omterm{def:ml:large-language-models-in-finance:prompt}{prompt} qui les contient (Lewis
et coauteurs, 2020). Le \emph{BM25}\index{BM25} est le score lexical de référence de la recherche~: pour chaque terme de la requête présent dans un
passage, l'inverse de la fréquence documentaire du terme, multipliée par une fonction saturante de son nombre d'occurrences dans le passage, normalisée par la
longueur du passage (Robertson et Zaragoza, 2009).
\end{definition}

La collection du chapitre compte 540 passages issus de 135 dépôts annuels synthétiques, un passage par entreprise, exercice et
indicateur (chiffre d'affaires, dette nette, effectifs, dépenses d'investissement), chacun rédigé avec l'un de trois synonymes («~revenue~», «~sales~»
ou «~turnover~») et deux mots de contexte. Une question est posée sur chaque passage, formulée avec un synonyme tiré au hasard, de sorte
qu'environ un tiers emploie le mot propre au passage~; la réponse est extraite du premier passage classé, ce qui fait qu'un bon passage
donne le bon chiffre. La \Cref{fig:ml:llm:retrieval} compare quatre recherches. Le \omterm{def:ml:large-language-models-in-finance:rag}{BM25} (\cref{lst:ml:llm:bm25}) retrouve chaque
passage visé avec ses propres mots et 5.0~\% des autres~: une question sur «~sales~» aboutit au passage d'une autre entreprise
La recherche dense, soit le cosinus entre les moyennes des \omterm{def:ml:text-from-bag-of-words-to-embeddings:embedding}{plongements de mots} de la question et du passage, sait que
«~sales~» veut dire chiffre d'affaires (les plongements, entraînés sur 5\,000 phrases générales d'affaires, placent les trois synonymes à un cosinus
supérieur à 0.99), mais pas quelle entreprise est visée~: 1.1~\%. Ajouter les deux scores standardisés n'aide guère (18.2~\%). Élargir
la requête par les plus proches voisins de ses mots dans l'espace des plongements avant de lancer \omterm{def:ml:large-language-models-in-finance:rag}{BM25} retrouve 89.1~\% des questions reformulées et 86.3~\% des
autres~: 88.1~\% au total.

\begin{omfigure}
\begin{tikzpicture}
\begin{axis}[width=10.5cm,height=5.8cm,ybar,bar width=12pt,ymin=0,ymax=110,ylabel={bon passage (\%)},area legend,
  xtick={0,1,2,3},xticklabels={BM25,dense,hybride,BM25 élargi},tick label style={font=\small},label style={font=\small},
  enlarge x limits=0.18,grid=major,grid style={gray!25},legend style={font=\footnotesize,at={(0.5,-0.2)},anchor=north,
  draw=none,fill=none,/tikz/every even column/.append style={column sep=8pt}},legend columns=2,table/col sep=comma]
\addplot[fill=omDef,draw=none] table[x=i,y=same]{figdata/ml/14-large-language-models-in-finance/retrieval.csv};
\addplot[fill=omThm,draw=none] table[x=i,y=other]{figdata/ml/14-large-language-models-in-finance/retrieval.csv};
\legend{la question reprend le mot du passage,la question emploie un synonyme}
\end{axis}
\end{tikzpicture}

\omcaption{Part des 540 questions dont le premier passage classé est celui visé, selon la méthode de recherche et la formulation. Données~:
\texttt{ml\_llm.retrieval}.}\label{fig:ml:llm:retrieval}
\end{omfigure}

\begin{definition}[Hallucination]\label{def:ml:large-language-models-in-finance:hallucination}
Une \emph{hallucination}\index{hallucination} est une réponse fluide qui n'est étayée ni par les données fournies au modèle ni par les faits~: un
chiffre, une citation ou une référence produits parce qu'il s'agit d'un texte probable, et non parce qu'il est vrai.
\end{definition}

Un système de recherche qui répond toujours hallucine par construction sur les questions auxquelles sa collection ne peut pas répondre. Sur 300
questions portant sur des couples entreprise-année sans dépôt, la recherche élargie renvoie à chaque fois un passage et un chiffre. Un contrôle
d'ancrage, qui ne répond que lorsque le passage nomme l'entreprise et l'exercice de la question, s'abstient sur les 300~; sur les questions
auxquelles il peut répondre, il refuse 11.9~\%, soit exactement les 64 sur 540 dont le premier passage concernait une autre entreprise ou un autre exercice, ce qui aurait
été une erreur. Un modèle génératif placé au-dessus de la recherche exige la même rigueur~: citer le passage, vérifier les entités et
s'abstenir plutôt que de deviner.

\section{Évaluation sans fuite depuis la date de coupure de l'entraînement}

\begin{definition}[Mémorisation, test d'anonymisation]\label{def:ml:large-language-models-in-finance:memo}
La \emph{mémorisation}\index{mémorisation} est la capacité d'un modèle à reproduire des exemples d'entraînement précis, y compris leurs
résultats, plutôt qu'une règle qui se généralise~; les \omterm{def:ml:large-language-models-in-finance:lm}{modèles de langage} mémorisent un texte vu ne serait-ce qu'une fois (Carlini et coauteurs, 2021).
Le \emph{test d'anonymisation}\index{test d'anonymisation} remplace les identifiants d'une entrée (noms d'entreprises, tickers, dates)
par des marqueurs et mesure de combien baisse la précision du modèle~: un modèle qui lit le texte perd peu, un modèle qui
se rappelle l'événement perd ce souvenir.
\end{definition}

Évaluer un \omterm{def:ml:large-language-models-in-finance:lm}{modèle de langage} commercial sur des titres historiques relève, sous une forme nouvelle, du biais d'anticipation (look-ahead bias, livre~7, chapitre~3)~: le
modèle peut avoir lu le titre et ce que l'action a fait ensuite. La \Cref{fig:ml:llm:contamination} mesure les deux tests sur les
deux modèles du chapitre. Sur les titres antérieurs au jour 907, que les deux ont vus avec leurs résultats, le modèle propre a raison dans 69.7~\% des
cas, et le modèle contaminé dans 65.4~\% avec les noms et les semaines, et dans 57.7~\% et 57.2~\% sans~: cet écart est de la mémoire. Entre les deux dates de coupure, le modèle contaminé a raison
dans 64.8~\% des cas, et le modèle propre dans 56.2~\%, et l'anonymisation ramène le modèle contaminé à 57.0~\% et laisse le
modèle propre là où il était (56.0~\%). Après les deux dates de coupure, le modèle contaminé a raison dans 59.2~\% des cas, le propre dans 56.3~\%, les deux
restant sous les 61.6~\% qu'atteint le signe de l'effet injecté.

\begin{omfigure}
\begin{tikzpicture}
\begin{axis}[width=11.5cm,height=6cm,ybar,bar width=8pt,ymin=50,ymax=72,ylabel={bonne direction (\%)},area legend,
  xtick={0,1,2},xticklabels={{avant les deux\\dates de coupure},{entre les dates\\de coupure},{après les deux\\dates de coupure}},xticklabel style={align=center},tick label style={font=\small},
  label style={font=\small},xmin=-0.5,xmax=2.5,grid=major,grid style={gray!25},legend columns=2,
  legend style={font=\footnotesize,at={(0.5,-0.3)},anchor=north,draw=none,fill=none,/tikz/every even column/.append style={column sep=8pt}},table/col sep=comma]
\addplot[fill=omDef,draw=none,bar shift=-13.5pt] table[x=i,y=clean]{figdata/ml/14-large-language-models-in-finance/contamination.csv};
\addplot[fill=omDef!40,draw=none,bar shift=-4.5pt] table[x=i,y=cleananon]{figdata/ml/14-large-language-models-in-finance/contamination.csv};
\addplot[fill=omThm,draw=none,bar shift=4.5pt] table[x=i,y=contam]{figdata/ml/14-large-language-models-in-finance/contamination.csv};
\addplot[fill=omThm!40,draw=none,bar shift=13.5pt] table[x=i,y=contamanon]{figdata/ml/14-large-language-models-in-finance/contamination.csv};
\addplot[black!70,dashed,thick,sharp plot,forget plot] table[x=x,y=y]{figdata/ml/14-large-language-models-in-finance/ceiling.csv};
\legend{propre,propre anonymisé,contaminé,contaminé anonymisé}
\end{axis}
\end{tikzpicture}

\omcaption{Part des titres dont chaque modèle prédit correctement la direction, par période, avec les noms et les semaines, et avec
des marqueurs. Les données d'entraînement du modèle propre s'arrêtent à la première date de coupure, celles du modèle contaminé à la seconde. En pointillés~: le
signe de l'effet injecté. Données~: \texttt{ml\_llm.contamination}.}\label{fig:ml:llm:contamination}
\end{omfigure}

Le résultat nommé du chapitre suit. La \emph{prime de \omterm{def:ml:large-language-models-in-finance:memo}{mémorisation}}, soit la précision entre les deux dates de coupure moins la précision
après les deux, vaut 5.6 points pour le modèle contaminé et $-0.1$ pour le modèle propre~; la \emph{perte due à l'anonymisation}, entre les
dates de coupure, vaut 7.8 points pour le modèle contaminé et 0.2 pour le modèle propre. Les deux tests détectent la contamination sans
connaître les données d'entraînement, et c'est précisément leur rôle~: un desk qui évalue un modèle dont il doute de la date de coupure les lance, et
ne fait confiance qu'aux chiffres postérieurs à la date de coupure annoncée. La littérature a observé les mêmes deux effets sur de vrais modèles. Lopez-Lira et
Tang (2023) ont évalué GPT-4 sur des titres publiés après sa date de coupure des connaissances, pour cette raison~; Glasserman et Lin (2023)
ont comparé des titres originaux à des titres débarrassés des identifiants de l'entreprise et ont constaté, dans l'échantillon, que les versions anonymisées
donnaient de meilleurs résultats~: la connaissance générale de l'entreprise distrayait davantage le modèle que le biais d'anticipation ne l'aidait. Le jouet n'a aucune
connaissance générale qui pourrait le distraire, et après les deux dates de coupure, l'anonymisation coûte à chaque modèle environ un point (59.2~\% à 58.1~\%,
56.3~\% à 55.3~\%)~: le \omterm{def:ml:large-language-models-in-finance:memo}{test d'anonymisation} est un test de mémoire, et sa perte doit se lire à l'aune de ce faible coût.

\begin{method}[Évaluer un modèle de langage sur des données de marché]\label{met:ml:llm:method}
\begin{enumerate}
\item Établir la \omterm{def:ml:large-language-models-in-finance:prompt}{date de coupure de l'entraînement} à partir de la documentation du fournisseur, et traiter tout ce qui la précède comme dans l'échantillon.
\item Ne calculer les scores que sur les textes publiés après la date de coupure~; si un historique antérieur doit être utilisé, présenter le \omterm{def:ml:large-language-models-in-finance:memo}{test d'anonymisation} et
la précision par période, et s'attendre à ce qu'ils divergent.
\item Figer la version du modèle, le \omterm{def:ml:large-language-models-in-finance:prompt}{prompt} et les paramètres de décodage dans le journal de recherche (livre~7, chapitre~1)~: une
mise à jour silencieuse du fournisseur constitue un nouveau modèle.
\item Comparer à une référence peu coûteuse (un dictionnaire, un modèle tf-idf) sur les mêmes titres et le même horizon.
\end{enumerate}
\end{method}

\section{Agents dans les flux de travail de recherche}

\begin{definition}[Agent à modèle de langage, appel d'outil]\label{def:ml:large-language-models-in-finance:agent}
Un \emph{agent à modèle de langage}\index{agent à modèle de langage} est une boucle dans laquelle un \omterm{def:ml:large-language-models-in-finance:lm}{modèle de langage} choisit des actions, observe leurs
résultats et décide de l'action suivante jusqu'à ce qu'une tâche soit achevée (Yao et coauteurs, 2023). Un \emph{appel d'outil}\index{appel d'outil} est
l'une de ces actions~: le modèle émet une requête structurée (un nom de fonction et des arguments) que le programme qui l'entoure exécute,
en renvoyant le résultat au modèle.
\end{definition}

Un agent capable de chercher, de calculer et de lancer des backtests est un assistant de recherche qui présente les mêmes modes de défaillance qu'un analyste junior,
à une vitesse de machine~: il utilisera des données qu'il ne devrait pas utiliser. La \Cref{lst:ml:llm:tools} est le harnais de la maison~: les outils sont enregistrés par
nom, un outil daté reçoit la date de référence de la tâche, et tout résultat publié après celle-ci est refusé, et chaque appel,
qu'il soit autorisé ou refusé, est consigné dans le journal. L'agent du chapitre est un plan fixe plutôt qu'un \omterm{def:ml:large-language-models-in-finance:lm}{modèle de langage}, de façon que son comportement
puisse être testé~: interrogé, à une date tirée au hasard, sur le dernier chiffre annuel d'une entreprise, il effectue une recherche, retient les passages qui concernent
l'entreprise et la métrique, et répond à partir du dernier exercice fiscal trouvé. Sur 400 questions de ce type, sans le garde-fou, 80.8~\%
des réponses proviennent de dépôts publiés après la date de référence et 19.2~\% sont exactes~; avec le garde-fou, aucune n'en provient et toutes sont exactes, et
chaque question laisse un appel journalisé.

\section{Ce qu'un desk peut et ne peut pas déléguer}

Les \omterm{def:ml:large-language-models-in-finance:lm}{modèles de langage} excellent dans le travail situé entre les données et la décision, qui coûte cher aux personnes et se vérifie à bon compte~: extraire
des champs dans les dépôts, résumer une conférence téléphonique en la comparant à la précédente, rédiger du code et des tests, étiqueter les actualités par événements et
entités, répondre à des questions sur une collection en citant le passage. Chacune de ces tâches produit un résultat vérifiable. Ce qui ne peut pas
être délégué, c'est ce qui ne se vérifie pas à bon compte~: la décision de trader, le jugement selon lequel un backtest est propre, l'affirmation
qu'un chiffre est juste alors qu'aucun passage ne la soutient. Trois règles fixent la frontière. Rien de ce que renvoie le modèle entre dans une
position ou dans un rapport sans une source qu'une personne ou un test peut vérifier. Les informations privilégiées (livre~7,
chapitre~12) restent à l'écart des \omterm{def:ml:large-language-models-in-finance:prompt}{prompts} envoyés à des services que la maison ne contrôle pas, puisqu'un \omterm{def:ml:large-language-models-in-finance:prompt}{prompt} est une divulgation. Et chaque appel de modèle
qui alimente la recherche est consigné avec sa version, son \omterm{def:ml:large-language-models-in-finance:prompt}{prompt} et sa sortie, de sorte que le résultat puisse être reproduit ou, à défaut,
expliqué. Kim, Muhn et Nikolaev (2024) rapportent que GPT-4, à qui l'on donnait des états financiers standardisés et anonymes, a prédit
la direction des variations de résultats mieux que les analystes~; c'est l'anonymisation qui a rendu le résultat crédible.

\section{Tutoriel~: le modèle qui savait}\label{tut:ml:large-language-models-in-finance:tut}

\begin{tutorial}
\textbf{Objectif.} Entraîner un \omterm{def:ml:large-language-models-in-finance:lm}{modèle de langage} propre et un modèle contaminé sur les titres du chapitre~13, et mesurer la prime de
\omterm{def:ml:large-language-models-in-finance:memo}{mémorisation} et la perte due à l'anonymisation~; comparer quatre méthodes de recherche et un contrôle d'ancrage~; lancer l'agent avec et sans son
garde-fou. \textbf{État final~:} \cref{fig:ml:llm:contamination,fig:ml:llm:retrieval}.
\begin{tutsteps}
\item \textbf{Un \omterm{def:ml:large-language-models-in-finance:lm}{modèle de langage} à décodeur seul, à l'échelle du jouet.}
\omcode{code/firm/llmeval/firm_llmeval.py}{51}{67}{Un modèle de langage causal à transformer.}\label{lst:ml:llm:tinylm}
\item \textbf{\omterm{def:ml:large-language-models-in-finance:rag}{BM25}.}
\omcode{code/firm/llmeval/firm_llmeval.py}{122}{143}{Le score BM25.}\label{lst:ml:llm:bm25}
\item \textbf{Outils point-in-time.}
\omcode{code/firm/llmeval/firm_llmeval.py}{168}{192}{Un registre d'outils avec un garde-fou point-in-time et un journal.}\label{lst:ml:llm:tools}
\item \textbf{Lancer} \texttt{ml\_llm.contamination()} (environ une minute sur un cœur), \texttt{verdict()}, \texttt{retrieval()},
\texttt{unanswerable()}, \texttt{agent\_run()} et \texttt{fig\_llm.py}.
\end{tutsteps}
\textbf{Pour aller plus loin.} Entraîner le modèle contaminé sur moins d'\omterm{def:ml:neural-networks-for-noisy-tabular-data:backprop}{époques} et observer la prime rétrécir~; remplacer les noms d'entreprises
entre les titres (un test de permutation d'entités) au lieu de les supprimer.
\end{tutorial}

\section{Construction~: évaluer les modèles de langage sans fuite}\label{bld:ml:large-language-models-in-finance:llmeval}

\begin{build}
\bfield{Rôle} \omterm{def:ml:large-language-models-in-finance:lm}{Modèles de langage} et recherche utilisés là où leur sortie peut être vérifiée, et évalués seulement là où leurs
données d'entraînement ne peuvent pas avoir vu la réponse.
\bfield{Interface} \texttt{Vocab}, \texttt{TinyLM(vocab\_size, d, heads, layers, max\_len)},
\texttt{train\_lm(model, seqs, epochs, lr, batch, seed)}, \texttt{next\_logprobs(model, prefixes)}~; \texttt{BM25(docs, k1, b)}, \texttt{dense\_scores},
\texttt{expand\_query}, \texttt{hybrid}~; \texttt{ToolRegistry(as\_of)} avec \texttt{register(name, fn, dated)},
\texttt{call(name, **kw)} et \texttt{log}.
\bfield{Règles} Aucun appel réseau et aucun poids téléchargé dans les tests~; scores présentés par période par rapport à la \omterm{def:ml:large-language-models-in-finance:prompt}{date de
coupure de l'entraînement}~; chaque \omterm{def:ml:large-language-models-in-finance:agent}{appel d'outil} journalisé~; les outils datés refusent tout élément postérieur à la date de référence.
\bfield{Tests d'acceptation} \texttt{code/firm/llmeval/tests/}~: le \omterm{def:ml:large-language-models-in-finance:lm}{modèle de langage} mémorise un petit corpus et est causal (la
prédiction d'un jeton ne change pas avec les jetons suivants)~; l'entraînement est déterministe~; \omterm{def:ml:large-language-models-in-finance:rag}{BM25} coïncide avec sa formule sur un exemple calculé à la main~;
l'expansion de requête n'ajoute que des voisins proches~; le registre refuse un document futur, journalise chaque appel et rejette les outils
inconnus.
\bfield{Pour aller plus loin} Un modèle à poids ouverts exécuté localement derrière la même interface, lancé hors ligne~; un test de permutation d'entités~; un registre de \omterm{def:ml:large-language-models-in-finance:prompt}{prompts} et de versions
dans l'outil de suivi d'expériences (chapitre~25).
\end{build}

\begin{omsources}
\item A.~Vaswani et coauteurs, «~Attention is all you need~», arXiv:1706.03762, 2017.
\item P.~Lewis et coauteurs, «~Retrieval-augmented generation for knowledge-intensive NLP tasks~», arXiv:2005.11401, 2020.
\item S.~Robertson et H.~Zaragoza, «~The probabilistic relevance framework: BM25 and beyond~», \emph{Foundations and Trends
in Information Retrieval} 3(4), 2009.
\item N.~Carlini et coauteurs, «~Extracting training data from large language models~», arXiv:2012.07805, 2021.
\item A.~Lopez-Lira et Y.~Tang, «~Can ChatGPT forecast stock price movements? Return predictability and large language
models~», arXiv:2304.07619, 2023.
\item P.~Glasserman et C.~Lin, «~Assessing look-ahead bias in stock return predictions generated by GPT sentiment
analysis~», arXiv:2309.17322, 2023.
\item S.~Wu et coauteurs, «~BloombergGPT: a large language model for finance~», arXiv:2303.17564, 2023.
\item S.~Yao et coauteurs, «~ReAct: synergizing reasoning and acting in language models~», arXiv:2210.03629, 2023.
\item A.~Kim, M.~Muhn et V.~Nikolaev, «~Financial statement analysis with large language models~», arXiv:2407.17866, 2024.
\end{omsources}

\section{Exercices}

\begin{exercise}[$\star$]\label{exo:ml:large-language-models-in-finance:1}
La documentation d'un modèle indique une \omterm{def:ml:large-language-models-in-finance:prompt}{date de coupure de l'entraînement} fixée au 30 juin. Un backtest couvre la période allant de janvier de l'année précédente à décembre.
Quelle partie de ce backtest permet de mesurer le jugement du modèle, et que mesure le reste~?
\end{exercise}

\begin{exercise}[$\star$]\label{exo:ml:large-language-models-in-finance:2}
Calculer la contribution \omterm{def:ml:large-language-models-in-finance:rag}{BM25} d'un terme de fréquence documentaire 45 parmi 540 passages, apparaissant une fois dans un passage de
longueur moyenne, avec $k_1 = 1.5$ et $b = 0.75$.
\end{exercise}

\begin{exercise}[$\star$]\label{exo:ml:large-language-models-in-finance:3}
À partir de la \cref{fig:ml:llm:contamination}, calculer la prime de \omterm{def:ml:large-language-models-in-finance:memo}{mémorisation} et la perte d'anonymisation de chaque modèle.
\end{exercise}

\begin{exercise}[$\star\star$]\label{exo:ml:large-language-models-in-finance:4}
Pourquoi la recherche dense seule ne retrouve-t-elle presque aucun des bons passages, et pourquoi ajouter son score à celui de \omterm{def:ml:large-language-models-in-finance:rag}{BM25} aide-t-il si
peu~?
\end{exercise}

\begin{exercise}[$\star\star$]\label{exo:ml:large-language-models-in-finance:5}
Pourquoi le modèle contaminé est-il meilleur que le modèle non contaminé après les deux dates de coupure (59.2~\% contre 56.3~\%)~? Est-ce une preuve de
jugement~?
\end{exercise}

\begin{exercise}[$\star\star$]\label{exo:ml:large-language-models-in-finance:6}
\emph{Trouver la faille.} «~Nous avons demandé au modèle le sens de 10\,000 titres de 2015 à 2023 et il a eu raison dans 70~\% des
cas, bien au-dessus des 56~\% de notre modèle tf-idf~; nous avons retiré les noms d'entreprises à titre de contrôle et il était encore à 65~\%.~»
\end{exercise}

\begin{exercise}[$\star\star\star$]\label{exo:ml:large-language-models-in-finance:7}
\emph{Programmation.} Vérifier que \texttt{TinyLM} est causal~: modifier le dernier jeton d'un \omterm{def:ml:large-language-models-in-finance:prompt}{prompt} et contrôler que les prédictions du
modèle à toutes les positions antérieures restent inchangées. Qu'est-ce qui, dans le code, le garantit~?
\end{exercise}

\begin{exercise}[$\star\star\star$]\label{exo:ml:large-language-models-in-finance:8}
Concevoir une tâche pour un agent (question de recherche, outils, date de référence, contrôles) pour «~lequel de nos signaux s'est le plus dégradé le trimestre dernier~?~»,
et lister ce que le journal doit contenir pour qu'un relecteur puisse reproduire la réponse.
\end{exercise}

\section{Problème~: Le modèle qui savait}

\begin{problem}[{Problème du week-end --- mémoire ou jugement}]\label{pb:ml:large-language-models-in-finance:1}
Les deux \omterm{def:ml:large-language-models-in-finance:lm}{modèles de langage} du chapitre, son corpus de dépôts réglementaires et son agent.

\medskip\noindent\textbf{Partie I --- Les modèles.}
\begin{enumerate}
\item Que voit chaque modèle à l'entraînement, et comment l'interroge-t-on sur un titre~?
\item Quelle est la taille du modèle et celle de son vocabulaire, et pourquoi le protocole compte-t-il davantage que la taille~?
\item Quelle est la meilleure exactitude disponible après les deux dates de coupure, et pourquoi n'est-elle pas plus élevée~?
\item Pourquoi un titre d'entraînement sur dix est-il anonymisé~?
\end{enumerate}

\medskip\noindent\textbf{Partie II --- Contamination.}
\begin{enumerate}[resume]
\item Quelles sont les exactitudes des deux modèles avant les deux dates de coupure, et que montre là l'exactitude anonymisée~?
\item Quelles sont-elles entre les deux dates de coupure, et qu'est-ce qui explique l'écart~?
\item Quel effet l'anonymisation a-t-elle sur chaque modèle entre les deux dates de coupure~?
\item Que dit la littérature de ces deux mêmes effets dans les modèles réels~?
\end{enumerate}

\medskip\noindent\textbf{Partie III --- Recherche et agents.}
\begin{enumerate}[resume]
\item Comment se comportent les quatre recherches avec la question formulée dans les termes du passage, puis avec d'autres mots~?
\item Que coûte le contrôle d'ancrage et que rapporte-t-il~?
\item Que fait l'agent sans son garde-fou~?
\item Que doit contenir le journal d'une exécution d'agent~?
\end{enumerate}

\medskip\noindent\textbf{Partie IV --- Le verdict.}
\begin{enumerate}[resume]
\item Énoncer le \emph{résultat nommé}~: la prime de \omterm{def:ml:large-language-models-in-finance:memo}{mémorisation} et la perte d'anonymisation pour le modèle contaminé et le
modèle non contaminé.
\item Comment évalue-t-on un modèle de sentiment sur l'actualité d'un fournisseur, construit sur un LLM~?
\item Quelles tâches confierait-on d'abord à un \omterm{def:ml:large-language-models-in-finance:lm}{modèle de langage} sur un desk de recherche~?
\item Qu'est-ce qui ne doit jamais entrer dans un \omterm{def:ml:large-language-models-in-finance:prompt}{prompt} envoyé hors de la société~?
\item Pourquoi le backtest d'un agent peut-il être moins bon que celui d'une personne~?
\item Comment rend-on reproductible un résultat de recherche fondé sur un LLM~?
\item Dans quel cas faut-il affiner un modèle plutôt que de le piloter par un \omterm{def:ml:large-language-models-in-finance:prompt}{prompt}~?
\item En une phrase~: que mesure l'exactitude d'un \omterm{def:ml:large-language-models-in-finance:lm}{modèle de langage} sur des données historiques~?
\end{enumerate}
\end{problem}

\section{Questions d'entretien}

\begin{interviewq}[$\star$ \iqroles{chercheur, ingénieur ML}]\label{iq:ml:large-language-models-in-finance:1}
Qu'est-ce qu'une \omterm{def:ml:large-language-models-in-finance:prompt}{date de coupure de l'entraînement}, et pourquoi compte-t-elle pour un backtest qui utilise un LLM~?
\end{interviewq}

\begin{interviewq}[$\star\star$ \iqroles{ingénieur ML}]\label{iq:ml:large-language-models-in-finance:2}
Pouvez-vous expliquer la \omterm{def:ml:large-language-models-in-finance:rag}{génération augmentée par la recherche}~? Comment évalueriez-vous l'étape de recherche séparément de la génération~?
\end{interviewq}

\begin{interviewq}[$\star\star$ \iqroles{chercheur}]\label{iq:ml:large-language-models-in-finance:3}
Comment testeriez-vous si la capacité apparente de prévision d'un LLM relève de la \omterm{def:ml:large-language-models-in-finance:memo}{mémorisation}~?
\end{interviewq}

\begin{interviewq}[$\star\star$ \iqroles{ingénieur ML, développeur}]\label{iq:ml:large-language-models-in-finance:4}
\omterm{def:ml:large-language-models-in-finance:rag}{BM25} ou plongements pour la recherche dans les dépôts réglementaires~? Pouvez-vous défendre un choix de conception~?
\end{interviewq}

\begin{interviewq}[$\star\star$ \iqroles{ingénieur ML}]\label{iq:ml:large-language-models-in-finance:5}
Comment empêcheriez-vous un agent de recherche d'utiliser des données postérieures à la date de référence, et comment prouveriez-vous qu'il ne l'a pas fait~?
\end{interviewq}

\begin{interviewq}[$\star\star\star$ \iqroles{chercheur, trader}]\label{iq:ml:large-language-models-in-finance:6}
Un fournisseur vend un score de sentiment fondé sur un LLM, affichant un ratio de Sharpe de 3 sur dix ans de backtest. Que demandez-vous avant un essai~?
\end{interviewq}
